\section{Technical estimates for normal coordinates}\label{sec:analytic-details}
We use the relative setting and notation of Section~\ref{sec:estimate}, with a bounded normal coordinate $Z$, selfadjoint coefficients annihilated by $\delta$, and actual vacuum $\xi$.
The vector $\xi_\mu$ is the marginal vacuum from Lemma~\ref{lem:scalar}.

\subsection{Graph approximation for the mixed energy}\label{subsec:mixed-energy-graphs}
We justify the mixed energy identity~\eqref{eq:mixed-energy} in the original derivative graphs.
The smooth phases $h_t(z,w):=(\bar z-\bar w)^2/(|z-w|^2+t)$ have modulus at most one.
Their first-leg divided differences, the second-leg divided differences and their conjugate-denominator versions have squared modulus at most $8Q_t$ from \eqref{eq:positive-soft}.
For example, writing $a=z-v$, $b=w-v$, the numerator identity
\[
 \bar a^2(|b|^2+t)-\bar b^2(|a|^2+t)
 =\bar a\bar b(\bar a b-\bar b a)+t(\bar a-\bar b)(\bar a+\bar b)
\]
gives this bound by separating the area term and the remaining positive terms of $Q_t$.
Its integral is at most $24\|\xi_\mu\|_2^2$ by the scalar vacuum calculation.
Approximate $h_t$ in $C^1$ on a containing real box by polynomial tensors.
Weak closedness of the joint graph of the four tensor-leg derivatives puts both phases in that original joint graph closure as $t\downarrow0$.
Convex combinations turn the weak graph convergence into strong graph convergence while keeping uniform bounds for the tensors.
Their limiting first divided difference is
\[
 \frac{(\bar z-\bar v)(w-v)-(\bar w-\bar v)(z-v)}
 {(z-v)(w-v)(z-w)},
\]
whose squared modulus is $\kappa(z,w,v)^2$.
The same approximants work in the mixed tensor-leg graphs.
Multiplication by each fixed core outer coefficient is continuous in these graphs; the Leibniz rule therefore supplies actual joint graph approximants to $\delta p$ and $\delta^\#p$ for every $p\in\A$.
Their partial traces are uniformly operator-norm bounded and converge in the first-derivative graphs, so their products with each fixed $L^2$ vacuum converge in $H$.
The actual adjoint formula for each real slice is $T_a^*p=p\zeta_a-T_ap$.

Let $A_a:=(\delta_a\otimes1)\delta_a-(1\otimes\delta_a)\delta_a$.
The cross terms in its Leibniz rule cancel.
The just established graphs justify contracting these actual iterates and pairing with a polynomial.
For the sum $A_x+A_y$, direct evaluation at $Z,Z^*$ gives, after tracing the middle leg, the tensors $(z-w)g/2$ and $(\bar z-\bar w)g/2$; all omitted coefficients are annihilated.
For instance, if $N:=\bar a b-\bar b a$, the relevant numerator at $Z$ is $-N^2/(\bar a|b|^2|a-b|^2)=a\kappa(a,b,0)^2$ before the factor $1/2$.
Equality on generators and the outer Leibniz rule identify the whole defect with one half of the weighted liberation.
The real-slice energy calculation used in Lemma~\ref{lem:kernel-energy} now proves \eqref{eq:mixed-energy} on analytic tests.
Positive kernel approximations dominated by $g$ pass to this identity by \eqref{eq:gap-g}, the finite word rule and the differentiated binomial series.

\subsection{The density-square estimate}\label{subsec:density-square-proof}
\begin{proof}[Proof of Lemma~\ref{lem:density-square}]
We reduce the estimate to three harmonic-analysis results, checking their hypotheses for an auxiliary measure.
Write $b(x)^2:=\iint\kappa(x,y,z)^2\dd\mu(y)\dd\mu(z)$, $f:=M_1\mu+b$ and $\sigma:=\mu/f$ on the full-measure set where $f$ is finite.
The positive lower bound for $M_1\mu$ on its compact support makes $\sigma$ finite.
For every disk, $\sigma(B(a,r))\le2r$, since $M_1\mu(x)\ge\mu(B(a,r))/(2r)$ for $x\in B(a,r)$.
For every Borel set $Q$,
\[
 \iiint_{Q^3}\kappa^2\dd\sigma^3\le3\sigma(Q),\qquad
 \int f^3\dd\sigma=\int f^2\dd\mu\le2G(\mu).
\]
For the first estimate, assign each triple to a point with smallest $f$ and bound the reciprocal product by that point's $f^{-3}$.

Tolsa's global density-square estimate under linear growth \cite[Theorem~19.1, arXiv version]{Tol17}, applied to $\sigma|_Q$, bounds its sharp square integral by $c\sigma(Q)$.
Smooth density differences are positive averages of sharp ones with weights $s(-\psi'(s))\,ds$, $1\le s\le2$, so the same testing bound holds for the smooth operator $U$ with values in $\mathcal K:=L^2((0,\infty),dr/r)$.
Its scalar kernel at scale $r$ is supported on $|x-y|\le4r$, is bounded by $c/r$, and has both first gradients bounded by $c/r^2$.
The indicator tests $b_Q=1_Q$, the linear growth of $\sigma$ and these size and second-variable regularity estimates verify \cite[Theorem~1.3]{LM17} with $m=\alpha=1$.
Thus $U:L^2(\sigma)\to L^2(\sigma;\mathcal K)$ is bounded uniformly.
The $\mathcal K$-valued kernel has size $c/|x-y|$ and difference bound $c|x-x'|/|x-y|^2$ when $|x-x'|\le|x-y|/2$, in both variables.
Linear growth and annular summation give both integral H\"ormander conditions; the component formula represents $U$ away from the support of the input.
These are precisely the additional hypotheses in \cite[Definition~2.1 and Theorem~2.7]{GM00}, which yield the strong $L^3(\sigma)\to L^3(\sigma;\mathcal K)$ bound.
Since $\mu=f\sigma$, H\"older's inequality gives
\[
 \int\|Uf\|_{\mathcal K}^2f\dd\sigma
 \le\|Uf\|_{L^3(\sigma;\mathcal K)}^2\|f\|_{L^3(\sigma)}
 \le c_\psi\int f^3\dd\sigma.
\]
This proves \eqref{eq:density-square} for the original measure.
\end{proof}
